%%
%% ACL 2025 LaTeX Template
%% Template for Annual Meeting of the Association for Computational Linguistics
%%

\documentclass[11pt]{article}

% ACL Packages
\usepackage[2025]{acl}  % ACL style file (needs to be downloaded separately)
\usepackage[T1]{fontenc}
\usepackage{times}
\usepackage{latexsym}
\usepackage{hyperref}
\usepackage{url}
\usepackage{booktabs}
\usepackage{amsfonts}
\usepackage{amsmath}
\usepackage{amssymb}
\usepackage{graphicx}
\usepackage{algorithm}
\usepackage{algorithmic}
\usepackage{xcolor}
\usepackage{multirow}
\usepackage{tikz}

% For hyperlinks
\hypersetup{
    colorlinks=true,
    linkcolor=blue,
    citecolor=blue,
    urlcolor=blue
}

% Math operators
\DeclareMathOperator*{\argmin}{arg\,min}
\DeclareMathOperator*{\argmax}{arg\,max}

% For linguistics notation
\newcommand{\B}[1]{\textbf{#1}}
\newcommand{\I}[1]{\textit{#1}}

% Title
\title{Your Paper Title Here}

% Authors
\author{First Author$^{1}$, Second Author$^{2}$, Third Author$^{1}$\\
\\
$^{1}$Department of Computer Science, University A\\
$^{2}$Department of Linguistics, University B\\
\\
\texttt{\{first, third\}@university-a.edu}, \texttt{second@university-b.edu}
}

\begin{document}

\maketitle

% Abstract
\begin{abstract}
Your abstract should concisely summarize: (1) the problem addressed,
(2) the proposed method/approach, (3) key results and findings, and
(4) the significance of the contributions. Keep it under 200 words for ACL.
\end{abstract}

% Main content
\section{Introduction}
\label{sec:introduction}

Introduce your NLP/CL problem and motivate your work:
\begin{itemize}
    \item What is the linguistic or computational challenge?
    \item Why is it important for the field?
    \item What limitations exist in current approaches?
\end{itemize}

State your contributions explicitly:

\textbf{Contributions.} Our main contributions are:
\begin{enumerate}
    \item We propose [method] for [task], which [key innovation]
    \item We create/release [dataset/resource] containing [statistics]
    \item We demonstrate [quantitative results] on [benchmarks]
\end{enumerate}

\section{Related Work}
\label{sec:related}

Survey related work in multiple subareas:

\paragraph{Task-specific Work.} Prior work on this task includes \cite{author2023paper}.

\paragraph{Method-related Work.} Related methodological approaches \cite{method2022}.

\paragraph{Datasets and Resources.} Existing resources for this task \cite{dataset2021}.

\paragraph{Our Novelty.} How our work differs from and improves upon existing approaches.

\section{Task Definition}
\label{sec:task}

Define the task formally. For NLP tasks, this typically includes:

\subsection{Problem Formulation}

Given an input sequence $X = (x_1, x_2, \ldots, x_n)$ where $x_i$ represents tokens,
the task is to predict an output $Y$ which could be:
\begin{itemize}
    \item A sequence $Y = (y_1, y_2, \ldots, y_m)$ for sequence-to-sequence tasks
    \item A label $y \in \mathcal{Y}$ for classification tasks
    \item A structure (tree, graph) for structured prediction
\end{itemize}

\subsection{Challenges}

Identify key challenges specific to this NLP task:
\begin{itemize}
    \item Challenge 1: e.g., long-range dependencies
    \item Challenge 2: e.g., data scarcity
    \item Challenge 3: e.g., linguistic ambiguity
\end{itemize}

\section{Method}
\label{sec:method}

\subsection{Overview}

Describe your approach at a high level.

\subsection{Model Architecture}

Detail the model architecture. For neural models:

\begin{equation}
    h_i = \text{Encoder}(x_i, h_{i-1})
\end{equation}

\begin{equation}
    P(y_t | y_{<t}, X) = \text{softmax}(W \cdot \text{Decoder}(h, y_{<t}))
\end{equation}

\subsection{Training Objective}

Define the training objective:

\begin{equation}
    \mathcal{L}(\theta) = -\sum_{(X,Y) \in \mathcal{D}} \log P_\theta(Y|X)
\end{equation}

\subsection{Key Components}

Describe key components of your approach:

\paragraph{Component 1.} Describe the first key component.

\paragraph{Component 2.} Describe the second key component.

\subsection{Inference}

Describe the inference/decoding procedure.

\section{Data}
\label{sec:data}

\subsection{Datasets}

Describe the datasets used:

\begin{table}[t]
    \centering
    \caption{Dataset statistics.}
    \label{tab:datasets}
    \begin{tabular}{lrrr}
        \toprule
        Dataset & Train & Dev & Test \\
        \midrule
        Dataset 1 & 50,000 & 5,000 & 5,000 \\
        Dataset 2 & 30,000 & 3,000 & 3,000 \\
        Dataset 3 & 100,000 & 10,000 & 10,000 \\
        \bottomrule
    \end{tabular}
\end{table}

\subsection{Preprocessing}

Describe preprocessing steps (tokenization, normalization, etc.).

\subsection{Annotation}

If you created a new dataset, describe the annotation process.

\section{Experiments}
\label{sec:experiments}

\subsection{Experimental Setup}

\paragraph{Baselines.} List baseline systems:
\begin{itemize}
    \item \textbf{Baseline 1} \cite{ref1}: Brief description
    \item \textbf{Baseline 2} \cite{ref2}: Brief description
\end{itemize}

\paragraph{Metrics.} Define evaluation metrics:
\begin{itemize}
    \item \textbf{BLEU}: For machine translation
    \item \textbf{F1}: For sequence labeling
    \item \textbf{Accuracy}: For classification
\end{itemize}

\paragraph{Implementation Details.} Hyperparameters, hardware, training time.

\subsection{Main Results}

\begin{table*}[t]
    \centering
    \caption{Main results on benchmark datasets. Best results in \textbf{bold}.}
    \label{tab:main_results}
    \begin{tabular}{lcccccc}
        \toprule
        & \multicolumn{3}{c}{\textbf{Dataset 1}} & \multicolumn{3}{c}{\textbf{Dataset 2}} \\
        \cmidrule(lr){2-4} \cmidrule(lr){5-7}
        \textbf{Model} & Acc & F1 & BLEU & Acc & F1 & BLEU \\
        \midrule
        Baseline 1 & 85.2 & 0.84 & 32.1 & 83.5 & 0.82 & 30.5 \\
        Baseline 2 & 87.3 & 0.86 & 34.5 & 85.2 & 0.84 & 32.8 \\
        \textbf{Ours} & \textbf{91.5} & \textbf{0.90} & \textbf{38.2} & \textbf{89.8} & \textbf{0.88} & \textbf{36.1} \\
        \bottomrule
    \end{tabular}
\end{table*}

\subsection{Analysis}

\paragraph{Ablation Study.} Analyze contribution of each component.

\paragraph{Error Analysis.} Discuss common error patterns.

\paragraph{Qualitative Examples.} Show example outputs.

\begin{table}[t]
    \centering
    \caption{Example outputs from our model.}
    \label{tab:examples}
    \begin{tabular}{lp{5cm}}
        \toprule
        \textbf{Input} & [Example input text] \\
        \midrule
        \textbf{Gold} & [Ground truth output] \\
        \textbf{Pred} & [Model prediction] \\
        \bottomrule
    \end{tabular}
\end{table}

\section{Discussion}
\label{sec:discussion}

\paragraph{Findings.} Summarize key insights from experiments.

\paragraph{Limitations.} Acknowledge limitations:
\begin{itemize}
    \item Limitation 1
    \item Limitation 2
\end{itemize}

\paragraph{Ethical Considerations.} Discuss any ethical implications for NLP work.

\section{Conclusion}
\label{sec:conclusion}

Summarize contributions and suggest future directions.

\section*{Acknowledgments}

Thank reviewers, annotators, and funding sources.

\section*{Limitations}

ACL requires a limitations section. Discuss:
\begin{itemize}
    \item Computational requirements
    \item Dataset/language coverage
    \item Edge cases where the method fails
\end{itemize}

% Bibliography
\bibliographystyle{acl_natbib}
\bibliography{references}

% Appendix
\appendix

\section{Additional Implementation Details}
\label{app:implementation}

Full hyperparameters, code availability, and additional details.

\section{Additional Examples}
\label{app:examples}

More example outputs and analysis.

\end{document}
